\section{Đa thức Lagrange nội suy [Nâng cao]}
\subsection{Phân tích \& Thiết kế (M0)}

\begin{itemize}

    \item \textbf{Input:}

    File \texttt{points.csv} chứa:
    \begin{itemize}
        \item Số nguyên $n$ ($n \le 8$)
        \item $n$ điểm phân biệt $(x_i, y_i)$
    \end{itemize}

    File \texttt{query.csv} chứa:
    \begin{itemize}
        \item Số nguyên $m$
        \item $m$ giá trị cần tính $P(x_q)$
    \end{itemize}

    \item \textbf{Xử lý:}

    \begin{itemize}

        \item \textbf{Case biên:}

        \begin{itemize}
            \item $n = 0 \rightarrow$ không nội suy được
            \item $n = 1 \rightarrow P(x)=y_0$
        \end{itemize}

        \item \textbf{Module 1 (Đọc points.csv):}

        \begin{itemize}
            \item Hàm \texttt{readPoints(path, x, y)}
            \item Kiểm tra mở file
        \end{itemize}

        \item \textbf{Module 2 (Đọc query.csv):}

        \begin{itemize}
            \item Hàm \texttt{readQueries(path, queries)}
        \end{itemize}

        \item \textbf{Module 3 (Tính Lagrange basis):}

        \begin{itemize}
            \item 
            \[
            L_i(x)=\prod_{j \ne i}\frac{x-x_j}{x_i-x_j}
            \]
        \end{itemize}

        \item \textbf{Module 4 (Tính đa thức nội suy):}

        \begin{itemize}
            \item 
            \[
            P(x)=\sum_{i=0}^{n-1} y_i \cdot L_i(x)
            \]
            \item Không cần tính hệ số đa thức
        \end{itemize}

        \item \textbf{Module 5 (Ghi kết quả):}

        \begin{itemize}
            \item Ghi file \texttt{result.csv}
            \item Format:
\begin{lstlisting}
x_q,P(x_q)
\end{lstlisting}
        \end{itemize}

    \end{itemize}

    \item \textbf{Output:}

\begin{lstlisting}
x_q,P(x_q)
\end{lstlisting}

(với \texttt{setprecision(6)})

\end{itemize}

\subsection{Mã nguồn C++11 (Tách module)}

\begin{lstlisting}
#include <iostream>
#include <vector>
#include <fstream>
#include <iomanip>

using namespace std;

bool readPoints(const string& path,
                vector<double>& x,
                vector<double>& y) {

    ifstream fin(path);

    if (!fin) return false;

    int n;
    fin >> n;

    x.resize(n);
    y.resize(n);

    for (int i = 0; i < n; ++i) {

        char comma;

        fin >> x[i] >> comma >> y[i];
    }

    return true;
}

bool readQueries(const string& path,
                 vector<double>& queries) {

    ifstream fin(path);

    if (!fin) return false;

    int m;
    fin >> m;

    queries.resize(m);

    for (int i = 0; i < m; ++i) {
        fin >> queries[i];
    }

    return true;
}

double evaluateLagrange(const vector<double>& x,
                        const vector<double>& y,
                        double xq) {

    int n = x.size();

    if (n == 0) return 0.0;

    double result = 0.0;

    for (int i = 0; i < n; ++i) {

        double Li = 1.0;

        for (int j = 0; j < n; ++j) {

            if (i == j) continue;

            Li *= (xq - x[j]) / (x[i] - x[j]);
        }

        result += y[i] * Li;
    }

    return result;
}

bool writeResults(const string& path,
                  const vector<double>& queries,
                  const vector<double>& values) {

    ofstream fout(path);

    if (!fout) return false;

    fout << fixed << setprecision(6);

    for (size_t i = 0; i < queries.size(); ++i) {

        fout << queries[i]
             << ","
             << values[i]
             << "\n";
    }

    return true;
}

int main() {

    vector<double> x, y;

    if (!readPoints("points.csv", x, y)) {

        cerr << "Loi mo file points!\n";
        return 1;
    }

    vector<double> queries;

    if (!readQueries("query.csv", queries)) {

        cerr << "Loi mo file query!\n";
        return 1;
    }

    vector<double> values(queries.size());

    for (size_t i = 0; i < queries.size(); ++i) {

        values[i] =
            evaluateLagrange(x, y, queries[i]);
    }

    if (!writeResults("result.csv",
                      queries,
                      values)) {

        cerr << "Loi mo file output!\n";
        return 1;
    }

    return 0;
}
\end{lstlisting}

\subsection{Kế hoạch kiểm thử (Test Cases)}

\textbf{Test Case 1 (Thông thường)}

\textbf{Input (points.csv):}

\begin{lstlisting}
3
0,1
1,3
2,2
\end{lstlisting}

\textbf{Input (query.csv):}

\begin{lstlisting}
3
0
1.5
3
\end{lstlisting}

\textbf{Expected Output (result.csv):}

\begin{lstlisting}
0.000000,1.000000
1.500000,2.875000
3.000000,-2.000000
\end{lstlisting}

\vspace{0.5cm}

\textbf{Test Case 2 (Case biên -- n=1)}

\textbf{Input (points.csv):}

\begin{lstlisting}
1
2,5
\end{lstlisting}

\textbf{Input (query.csv):}

\begin{lstlisting}
2
0
10
\end{lstlisting}

\textbf{Expected Output (result.csv):}

\begin{lstlisting}
0.000000,5.000000
10.000000,5.000000
\end{lstlisting}

\vspace{0.5cm}

\textbf{Test Case 3 (Case biên -- n=0)}

\textbf{Input (points.csv):}

\begin{lstlisting}
0
\end{lstlisting}

\textbf{Input (query.csv):}

\begin{lstlisting}
2
1
2
\end{lstlisting}

\textbf{Expected Output (result.csv):}

\begin{lstlisting}
1.000000,0.000000
2.000000,0.000000
\end{lstlisting}